% ============================================================================
%  PAPER 2 - hardware follow-up to the sim-to-sim study. INTRO + METHODS ONLY.
%
%  FORMAT: IEEE RA-L. Per RA-L Information for Authors, INITIAL submissions use
%  the IEEE CONFERENCE format (ieeeconf.cls, from the RAS PaperCept support
%  files); the journal IEEEtran format is only for the accepted final version.
%    - US Letter, two column, 10pt
%    - 6 pages free, up to 2 extra at a per-page charge, references included
%    - DOUBLE-ANONYMOUS review: no author names, no identifying URLs, no
%      "our previous work" phrasing. Violations are desk-rejected.
%    - First two keywords must come from the official RA-L keyword list.
%
%  OVERLEAF: needs ieeeconf.cls in the same project. Overleaf does NOT ship it;
%  copy the one committed alongside this file.
%
%  \anonymoustrue for submission; \anonymousfalse for your own reading copy.
%  Placeholders print in red via \todo{}. Search "TODO" before submitting.
%  Sections V (Results) and VI (Discussion) are deliberate stubs.
% ============================================================================
\documentclass[letterpaper, 10pt, conference]{ieeeconf}
\IEEEoverridecommandlockouts

\usepackage{graphicx}
\usepackage{amsmath,amssymb}
\usepackage{booktabs}
\usepackage[dvipsnames]{xcolor}
\usepackage[colorlinks=true,allcolors=blue]{hyperref}

% ieeeconf descends from IEEEtran v1.6, which predates \IEEEPARstart
\providecommand{\IEEEPARstart}[2]{\PARstart{#1}{#2}}

\newif\ifanonymous
\anonymoustrue          % RA-L initial submission is double-anonymous

% Visible placeholder marker - delete the body of this macro before submission
\newcommand{\todo}[1]{\textcolor{red}{\textbf{[TODO: #1]}}}
\newcommand{\ie}{i.e.}
\newcommand{\eg}{e.g.}

\begin{document}

\title{\LARGE \bf Does Sim-to-Sim Predict Sim-to-Real? Multimodal Diffusion
Inverse Kinematics on a Physical Tendon-Driven Continuum Arm}

\ifanonymous
\author{Anonymous submission}
\else
\author{Adhvay Jagadeesh$^{1}$
\thanks{$^{1}$Independent Researcher. \todo{contact email}}}
\fi

\maketitle
\thispagestyle{empty}
\pagestyle{empty}

\begin{abstract}
Conditional diffusion models can represent the full, multimodal solution set
of a redundant inverse kinematics (IK) problem, and a prior controlled study
established three things in simulation: the advantage over deterministic
regressors is specifically a redundancy effect, the learned inverse map does
not survive a change of dynamics model, and the recovery mechanisms available
after that collapse are sharply bounded. That study was simulation-only. We
port its full protocol to a physical two-section, six-tendon continuum arm
with camera-based shape measurement, and ask two questions. First, does the
transfer collapse reproduce against real hardware, where the source model is
a calibrated model of the very arm being commanded? Second, and more useful
to the field, does the cheap sim-to-sim characterization \emph{rank} methods
the same way expensive hardware evaluation does---\ie, can sim-to-sim serve
as a screening proxy? We describe the platform, the measurement chain, and a
protocol in which every reported gap is bounded below by an independently
measured noise floor.
\todo{one sentence of headline results once measured}
\end{abstract}

\begin{keywords}
Modeling, Control, and Learning for Soft Robots; Deep Learning Methods;
inverse kinematics; sim-to-real transfer.
\end{keywords}

% ============================================================================
\section{Introduction}
% ============================================================================

\IEEEPARstart{A}{} tendon-driven continuum arm with three tendons per section
has more actuation degrees of freedom than task degrees of freedom: a
two-section arm exposes six actuator commands to specify a three-dimensional
tip position. The inverse map is therefore one-to-many, and its preimages are
not merely a connected null space---physically distinct arm shapes reach the
same point. A deterministic learned regressor trained on motor-babbling
data~\cite{giorelli2015neural} fits the conditional \emph{mean} of that
multimodal distribution and can return an actuation realizing none of the
true solutions, while a conditional diffusion
model~\cite{ho2020denoising,song2021denoising} learns the full conditional
distribution and can sample every solution mode.

A recent controlled simulation study~\cite{prior} established three results
about this setting. (i) The diffusion advantage is \emph{specifically}
redundancy resolution rather than generic model capacity: removing redundancy
by conditioning on whole-body shape collapses sampled diversity by roughly
$80\times$ and nearly erases the gap to a deterministic baseline, while
adding redundancy by increasing segment count widens it monotonically. (ii)
The learned inverse map is simulator-specific: actuations from a model
trained on a piecewise-constant-curvature (PCC) simulator collapse when
executed on a Cosserat-rod simulator of the same arm, on every training seed.
(iii) The recovery available after that collapse is bounded---candidate
selection under a target-domain query budget helps but never solves a target,
a diversity-aware query ordering significantly outperforms confidence
ordering, guidance-based adaptation is at best suggestive, gain domain
randomization achieves nothing, and target-domain fine-tuning buys genuine
transfer only across a sharp retention--transfer cliff that data mixing
cannot smooth.

Every one of those results was obtained between two simulators. That is a
deliberate methodological choice---it isolates model mismatch from
sensing noise, actuator imperfection, and unmodeled contact---but it leaves
the practically decisive question untouched. This paper addresses it
directly.

\subsection*{Why hardware changes the question}

Moving to a physical arm does more than add realism; it inverts the
relationship between the source model and the target domain. In the prior
study, the source simulator (PCC) was never intended as a model of the target
simulator (Cosserat rod); the two were independently specified, and their
disagreement was structural and directional. Here, by contrast, the source
model is a PCC model whose parameters are \emph{fitted to the very arm the
policy will command}. If the transfer collapse survives even that
best-case pairing, the finding is considerably stronger than its sim-to-sim
predecessor. If it does not, the magnitude of the sim-to-sim gap was an
artifact of pairing two unrelated models, and that is worth reporting with
equal prominence.

Hardware also makes one of that study's framings natural rather than
contrived. The query-budget experiment asked how many \emph{expensive}
target-domain evaluations candidate selection requires; in simulation an
``expensive'' evaluation was four seconds of Cosserat integration. On
hardware it is a physical actuation, a settling interval, and a camera
measurement---genuinely costly, non-batchable, and subject to wear. The
economics the experiment was designed to probe are real here.

\subsection*{Contributions}

\begin{enumerate}
  \item \textbf{An inexpensive, reproducible hardware platform and protocol.}
  A two-section, six-tendon continuum arm built from commodity parts
  (\todo{final bill-of-materials cost}), with camera-based shape measurement,
  and a full port of that study's evaluation suite. All fabrication
  templates, calibration procedures, and code are released.

  \item \textbf{A sim-to-real characterization of multimodal diffusion IK.}
  We measure whether the transfer collapse, the redundancy-specific
  advantage, and the bounded-recovery result each reproduce against physical
  hardware (Sec.~\ref{sec:results}).

  \item \textbf{A cross-domain consistency test.} We compare the
  \emph{ranking} of methods and recovery mechanisms under the prior
  sim-to-sim evaluation against their ranking on hardware. Agreement would
  license sim-to-sim screening as a cheap proxy for hardware evaluation;
  disagreement would caution against a practice the field currently relies on
  implicitly.
\end{enumerate}

Because a physical platform introduces error sources that simulation does
not, we treat measurement rigour as a first-class part of the method rather
than an appendix. Sec.~\ref{sec:noise} specifies an independently measured
noise floor, and no reported gap is interpreted unless it exceeds that floor.

% ============================================================================
\section{Related Work}
% ============================================================================

\textbf{Continuum-arm IK.} Non-learned approaches invert a model at query
time---Jacobian schemes for nonconstant-curvature
arms~\cite{giorelli2015neural} and real-time FEM
inversion~\cite{duriez2013control}. The constant-curvature kinematics
underlying our source model are reviewed in~\cite{webster2010design}, and
tendon-driven construction with spacer discs follows standard practice in the
continuum-robotics literature. Learned approaches fit a deterministic
target-to-actuation map on babbling data~\cite{giorelli2015neural}; mixture
density networks~\cite{bishop1994mixture} are the classical multimodal
alternative.

\textbf{Diffusion models for robot action.} Denoising diffusion
models~\cite{ho2020denoising} with DDIM sampling~\cite{song2021denoising} and
classifier-free guidance~\cite{ho2022classifier} underpin diffusion
policies~\cite{chi2023diffusion}. IKDiffuser~\cite{zhang2025ikdiffuser} solves
IK for rigid multi-arm systems with flexible kinematic-tree topology;
compliant continuum bodies are outside its scope.

\textbf{Crossing the reality gap.} Domain
randomization~\cite{tobin2017domain} and target-domain fine-tuning are the
standard levers, with catastrophic forgetting under naive fine-tuning
well-documented~\cite{kirkpatrick2017overcoming}. What is comparatively rare,
and what this paper supplies, is a single method evaluated under \emph{both}
a controlled sim-to-sim gap and a physical sim-to-real gap with a shared
protocol, permitting the two to be compared rather than merely reported
side by side.

% ============================================================================
\section{Hardware Platform}
\label{sec:hardware}
% ============================================================================

\subsection{Mechanical design}

The arm is a two-section tendon-driven continuum manipulator built around a
single flexible backbone. The backbone is a $4.76\,$mm ($3/16''$) structural
fiberglass rod of free length $305\,$mm ($12''$). Six spacer discs of
$38.1\,$mm ($1.5''$) diameter and $3.2\,$mm ($1/8''$) thickness are bonded to
the backbone at a uniform $50.8\,$mm ($2''$) pitch with two-part epoxy,
dividing the arm into two actuated sections of three discs each.

Each disc carries seven holes: a central clearance hole for the backbone and
six tendon guides arranged on two concentric circles at radii $8.1\,$mm and
$14.0\,$mm, with three holes per circle at $0^\circ$, $120^\circ$ and
$240^\circ$. Section~A's three tendons run on the inner circle and terminate
at disc~3; section~B's run on the outer circle, pass freely through discs 1--5,
and terminate at disc~6. Both sections therefore actuate along
\emph{identical} angular headings and differ only in moment arm---a
deliberate choice that preserves the three-fold symmetry assumed by the
constant-curvature source model, so that the two sections are related by a
single scalar per section rather than by a change of basis.

Tendons are $0.79\,$mm ($1/32''$) $7\times7$ stainless wire rope, terminated
with swaged sleeves and pre-tensioned through inline turnbuckles at the base.
Steel rope is used rather than braided polymer line specifically to suppress
creep: commanded tendon length must remain a stable function of actuator
angle across thousands of cycles.

All six discs are drilled through a single thickened master jig rather than
individually marked. This converts per-disc random error, which no
calibration can absorb, into a single systematic bias, which one gain fit per
section absorbs exactly. \todo{report measured disc-to-disc hole scatter}

\subsection{Actuation}

Six ROBOTIS DYNAMIXEL XL330-M288-T smart servos, one per tendon, are mounted
on a rigid $305\times305\,$mm acrylic base plate in three pairs at $120^\circ$
about the backbone, aligning each servo pair with the tendon heading it
drives. Each servo carries a metal horn with a brass standoff acting as a
cable capstan; tendon displacement is related to commanded servo angle by
$\Delta \ell = \theta\,(r_{s} + r_{c})$ with $r_{s}$ the capstan radius and
$r_{c}$ the cable radius, measured empirically rather than assumed
(Sec.~\ref{sec:calib}). Cable is wound in a single layer so that the
effective radius, and hence this relation, remains constant over the working
range. Servos operate in extended-position (multi-turn) mode with software
limits set below the tendon's mechanical travel.

The base plate is drilled with the same seven-hole pattern as a spacer disc
and functions as disc~V$_0$: tendons emerge from the plate on the same radii
and headings they occupy at every disc above, so that no tendon runs obliquely
between the base and the first disc.

\subsection{Command integrity}

One implementation detail is reported because it is invisible until it
corrupts data. The six servos share a single half-duplex bus and are
commanded with a synchronised broadcast write, so that both sections actuate
simultaneously rather than in sequence. That broadcast returns no
acknowledgement by design; the host library therefore reports success as soon
as the bytes are queued, whether or not any servo acted on them. If the next
packet is placed on the bus too soon, the servos discard the write
\emph{silently}---no error code, no fault flag, and every register still
reading back the expected configuration.

We characterised the threshold directly by varying the interval between the
broadcast write and the following read, and counting how many servos executed
the command: none at $0$ or $5\,$ms, and all six at $10\,$ms and above. All
experiments enforce a $25\,$ms interval, roughly $2.5\times$ the measured
threshold. Without it the failure does not announce itself---it appears in the
data as an actuation that was commanded but never applied, which is
indistinguishable from a genuine tracking failure and would inflate any
reported sim-to-real gap. We note it because anyone reproducing this platform
on a shared servo bus will meet it.

\subsection{Actuation parameterization}

To keep the interface identical to the prior study, the policy commands a
normalized actuation vector $\mathbf{q}\in[0,1]^{6}$, mapped affinely onto
each servo's calibrated travel range. Per-section curvature components follow
the same three-chamber form used for the pneumatic arm
in~\cite{prior}, with tendon displacements substituted for chamber pressures:
\begin{align}
k_x &= g\,(\ell_0 - \tfrac{1}{2}\ell_1 - \tfrac{1}{2}\ell_2), \\
k_y &= g\,\tfrac{\sqrt{3}}{2}\,(\ell_1 - \ell_2),
\end{align}
with $g$ a per-section gain fitted in Sec.~\ref{sec:calib}. This substitution
is what allows the entire software stack, the trained policies, and the
evaluation suite to transfer unmodified.

% ============================================================================
\section{Measurement System}
\label{sec:sensing}
% ============================================================================

\subsection{Vision chain}

Arm shape is measured by a fixed monocular global-shutter camera
(OV9281, $1280\times800$, monochrome) observing planar fiducial markers
mounted flat on the spacer discs. A global shutter is used rather than a
rolling shutter so that measurements taken during residual motion are not
geometrically distorted---relevant because, as in the prior study's target
domain, the arm is sampled at a fixed settling horizon rather than at true
static equilibrium.

Markers are ArUco fiducials~\cite{garrido2014automatic} of side
\todo{final marker size} mm, the largest that fits a $38.1\,$mm disc with an
adequate quiet zone. Markers on discs 3 and 6 give section-boundary and tip
pose; a full set on all six discs gives whole-backbone shape for the
shape-conditioning experiment. Detection and pose estimation use
OpenCV~\cite{bradski2000opencv}.

Camera intrinsics and distortion are estimated once from
\todo{$N$} checkerboard views and held fixed. The camera-to-base extrinsic
transform is established from a marker at a surveyed position on the base
plate, and the camera is not moved thereafter; every reported position is
expressed in the base frame.

\subsection{Sampling protocol}

One sample comprises: draw an actuation, command all six servos, wait a fixed
settling interval of \todo{settle time} s, capture one frame, detect markers,
and record the actuation together with the resulting disc poses. Detections
failing a confidence or visibility check are discarded and redrawn; we report
the discard rate as a data-quality statistic, a diagnostic with no analogue in
simulation.

The settling interval is fixed rather than adaptive, making the target domain
a deterministic dynamic-snapshot protocol---the same convention adopted, and
explicitly flagged, in that study's Cosserat domain. It is applied
identically to every sample, every dataset, and every evaluation here.

\subsection{Noise floor and repeatability}
\label{sec:noise}

Because the central claim concerns the size of a measured gap, we bound
below what a gap can mean. Two quantities are measured before any policy is
evaluated.

\emph{Sensing noise.} Two markers are fixed at a separation measured
independently with digital calipers, and the vision chain estimates that
separation \todo{$N$} times under the operating configuration. The spread of
those estimates is the sensing noise floor.
\todo{report mean error and standard deviation}

\emph{Actuation repeatability.} A set of \todo{$N$} actuations is commanded
\todo{$k$} times each, in randomized order with a full return to the home
configuration between repeats. The dispersion of measured tip position for a
repeated command is the actuation repeatability floor, and it folds in servo
backlash, cable hysteresis, and any residual creep.
\todo{report repeatability}

A reported transfer gap is interpreted only where it exceeds the combination
of these two floors. This is the substitute, on hardware, for the exact
determinism a simulator provides for free.

% ============================================================================
\section{Source Model, Training, and Protocol}
\label{sec:protocol}
% ============================================================================

\subsection{Fitted source model}
\label{sec:calib}

The source domain is a differentiable PCC model of \emph{this} arm, with
parameters fitted to it rather than assumed. Calibration proceeds in two
stages. First, each capstan's effective radius is measured directly by
commanding a known servo rotation and measuring paid-out cable with calipers,
yielding the angle-to-tendon-length relation. Second, the per-section
curvature gain $g$ is fitted by least squares over \todo{$N$} calibration
actuations, minimizing measured-versus-predicted disc pose.
\todo{report fitted gains and calibration residual}

This is the pairing that makes the hardware experiment sharp. The source
model is not an arbitrary alternative dynamics---it is the best
constant-curvature description of the physical arm that the data supports,
and its residual is reported so readers can judge how good a source model the
policy was given.

\subsection{Fit quality as a reported quantity}

Because the source model is fitted rather than assumed, its residual is a
result, not a footnote: it upper-bounds how much of any observed transfer gap
is attributable to an imperfect source model rather than to a genuine
limitation of the learned inverse map. We report the distribution of
measured-versus-predicted disc position error, and compare it against the
sensing noise floor of Sec.~\ref{sec:noise}. A residual close to that floor
means the constant-curvature description fits this arm about as well as the
instrumentation can resolve, and the transfer gap may be read at face value; a
residual far above it bounds how much of the gap is our own modelling error.
\todo{report fitted gains and residual distribution}

\subsection{Training data and models}

Training data are \todo{$N$} actuations drawn i.i.d.\ uniformly and pushed
through the fitted PCC model. We evaluate the same four methods as the prior
study, with architectures and hyperparameters carried over unchanged so that
any difference in outcome is attributable to the domain rather than to
retuning: a conditional diffusion model (residual FiLM-conditioned MLP
denoiser, DDPM training with cosine schedule and classifier-free guidance,
DDIM sampling), a deterministic MLP regressor, a mixture density
network~\cite{bishop1994mixture}, and a random-restart gradient-IK solver
optimizing through the differentiable fitted model. Unless stated otherwise
$K$ candidates are sampled per target and best-of-$K$ is reported; success is
tip error below \todo{tolerance} mm, set to the same fraction of arm length
used previously.

\subsection{Experiments}

\begin{itemize}
  \item \textbf{H1 (source-domain accuracy).} Tip error against the fitted
  PCC model, verifying the policies are trained correctly before hardware is
  involved. This isolates training failures from transfer failures.

  \item \textbf{H2 (multimodality on hardware).} Ground-truth solution sets
  are enumerated on the fitted source model, and mode recall is scored using
  \emph{measured} configurations. Distances are computed in curvature space
  rather than actuation space, since the tendon-to-curvature map carries a
  per-section null space that would otherwise contaminate distance-based
  metrics.

  \item \textbf{H3 (downstream value).} Obstacle avoidance with physical
  obstacles at surveyed positions; success requires some candidate to be both
  accurate and collision-free.

  \item \textbf{H4 (sim-to-real transfer).} The headline experiment:
  actuations from the source-trained policies executed on the arm, measured
  by camera, and compared against source-domain performance.

  \item \textbf{H5 (recovery mechanisms).} Candidate selection under a
  physical query budget, diversity-aware versus confidence query ordering,
  and fine-tuning on measured data across a range of source-data mixing
  ratios---testing in particular whether the retention--transfer cliff
  persists when the target-domain data are real measurements rather than
  simulator output.
\end{itemize}

\subsection{Cross-domain comparison}

The consistency test of Contribution~3 compares method and mechanism
\emph{rankings} between the prior sim-to-sim study and H1--H5 here. We
deliberately do not fit a relationship between the two gap magnitudes: with a
single hardware platform, that would be an $n{=}1$ regression. Rank agreement
is what a practitioner would actually rely on when using a cheap sim-to-sim
study to decide which method to carry forward to hardware, and rank agreement
is what we report.

% ============================================================================
\section{Results}
\label{sec:results}
% ============================================================================

\todo{Entire section. Populate from H1--H5. Suggested structure: Table I
comparing all four methods in source domain and on hardware; Fig.~2 query
budget curve with physical queries; Table II fine-tuning mixing grid;
Fig.~3 rank comparison against the sim-to-sim study.}

% ============================================================================
\section{Discussion and Limitations}
% ============================================================================

\todo{Entire section. Limitations to state explicitly once results exist:
single hardware platform ($n{=}1$ for the cross-domain claim); planar or
near-planar sensing if a single camera is used; open-loop IK formulation with
best-of-$K$ presupposing a verifier; fitted PCC as source model is a
best-case pairing and a harsher source model would likely widen the gap.}

% ============================================================================
\section*{Reproducibility}
% ============================================================================

\ifanonymous
Code, fabrication templates, calibration procedures, and per-experiment result
files are available in an anonymized repository:
\todo{anonymous.4open.science mirror URL - mint it from a fresh branch}.
\else
Code, fabrication templates, calibration procedures, and per-experiment result
files are available at \todo{public repository URL}.
\fi

% ============================================================================
\begin{thebibliography}{99}

\bibitem{prior}
\todo{The sim-to-sim paper. Under double-anonymous review cite it in the
third person as an existing preprint - never as "our previous work", and do
not include a repository URL that identifies the authors.}

\bibitem{ho2020denoising}
J.~Ho, A.~Jain, and P.~Abbeel, ``Denoising diffusion probabilistic models,''
in \emph{Advances in Neural Information Processing Systems}, 2020.

\bibitem{song2021denoising}
J.~Song, C.~Meng, and S.~Ermon, ``Denoising diffusion implicit models,'' in
\emph{International Conference on Learning Representations}, 2021.

\bibitem{ho2022classifier}
J.~Ho and T.~Salimans, ``Classifier-free diffusion guidance,'' \emph{arXiv
preprint arXiv:2207.12598}, 2022.

\bibitem{chi2023diffusion}
C.~Chi, S.~Feng, Y.~Du, Z.~Xu, E.~Cousineau, B.~Burchfiel, and S.~Song,
``Diffusion policy: Visuomotor policy learning via action diffusion,'' in
\emph{Robotics: Science and Systems}, 2023.

\bibitem{zhang2025ikdiffuser}
Z.~Zhang \emph{et al.}, ``IKDiffuser: A diffusion-based generative inverse
kinematics solver for kinematic trees,'' \emph{arXiv preprint
arXiv:2506.13087}, 2025.

\bibitem{giorelli2015neural}
M.~Giorelli, F.~Renda, M.~Calisti, A.~Arienti, G.~Ferri, and C.~Laschi,
``Neural network and Jacobian method for solving the inverse statics of a
cable-driven soft arm with nonconstant curvature,'' \emph{IEEE Transactions
on Robotics}, vol.~31, no.~4, pp. 823--834, 2015.

\bibitem{webster2010design}
R.~J.~Webster III and B.~A.~Jones, ``Design and kinematic modeling of constant
curvature continuum robots: A review,'' \emph{The International Journal of
Robotics Research}, vol.~29, no.~13, pp. 1661--1683, 2010.

\bibitem{duriez2013control}
C.~Duriez, ``Control of elastic soft robots based on real-time finite element
method,'' in \emph{IEEE International Conference on Robotics and Automation},
2013.

\bibitem{bishop1994mixture}
C.~M.~Bishop, ``Mixture density networks,'' Aston University, Tech. Rep.
NCRG/94/004, 1994.

\bibitem{tobin2017domain}
J.~Tobin, R.~Fong, A.~Ray, J.~Schneider, W.~Zaremba, and P.~Abbeel, ``Domain
randomization for transferring deep neural networks from simulation to the
real world,'' in \emph{IEEE/RSJ International Conference on Intelligent
Robots and Systems}, 2017.

\bibitem{kirkpatrick2017overcoming}
J.~Kirkpatrick \emph{et al.}, ``Overcoming catastrophic forgetting in neural
networks,'' \emph{Proceedings of the National Academy of Sciences}, vol.~114,
no.~13, pp. 3521--3526, 2017.

\bibitem{garrido2014automatic}
S.~Garrido-Jurado, R.~Mu\~{n}oz-Salinas, F.~J.~Madrid-Cuevas, and
M.~J.~Mar\'{i}n-Jim\'{e}nez, ``Automatic generation and detection of highly
reliable fiducial markers under occlusion,'' \emph{Pattern Recognition},
vol.~47, no.~6, pp. 2280--2292, 2014.

\bibitem{bradski2000opencv}
G.~Bradski, ``The OpenCV library,'' \emph{Dr. Dobb's Journal of Software
Tools}, 2000.

\end{thebibliography}

\end{document}
